\documentclass[11pt,a4paper]{article}
\usepackage{amsmath,amssymb}

\begin{document}
	
	\subsection*{Common Derivative Notations in Differential Equations}
	
	For an unknown function $y = y(x)$, the derivative with respect to $x$ can be written in several equivalent ways:
	
	\begin{itemize}
		\item Leibniz notation:
		$$
		\frac{dy}{dx}
		$$
		emphasizes the rate of change of $y$ with respect to $x$.
		
		\item Lagrange notation:
		$$
		y'
		$$
		is a compact form widely used in ODEs.
		
		\item Newton notation (common in physics, when $x$ represents time $t$):
		$$
		\dot{y}
		$$
		denotes $\dfrac{dy}{dt}$.
	\end{itemize}
	
	\paragraph{Example.}  
	The differential equation
	$$
	y' = f(x,y)
	$$
	is equivalently written as
	$$
	\frac{dy}{dx} = f(x,y),
	$$
	or, if $x=t$ represents time,
	$$
	\dot{y} = f(t,y).
	$$
	
	All three forms describe the same relationship: a rule for the derivative of the unknown function $y$ in terms of the variables $(x,y)$.
	
	
	\subsection*{Normal Notation for a Differential Equation}
	
	In common notation a first--order ODE is written as
	$$
	y' = f(x,y).
	$$
	
	Formally:
	\begin{itemize}
		\item $y = y(x)$ is the unknown function (solution).
		\item $y' = \dfrac{dy}{dx}$ is its derivative with respect to the independent variable $x$.
		\item $f : D \subseteq \mathbb{R}^2 \;\longrightarrow\; \mathbb{R}$ in the real case, or 
		$$f : D \subseteq \mathbb{R} \times \mathbb{C} \;\longrightarrow\; \mathbb{C}$$
		in the complex case.
		\item $D$ is the domain of $f$, i.e.\ the set of all pairs $(x,y)$ where $f(x,y)$ is defined.
	\end{itemize}
	
	Thus the equation
	$$
	y' = f(x,y)
	$$
	means: \emph{find a function $y : I \to \mathbb{C}$ such that for every $x \in I$, the derivative $y'(x)$ equals $f(x,y(x))$, where $(x,y(x)) \in D$.}
	
	
	\section*{Complex--Valued Functions and ODE Notation}
	
	\subsection*{Pedantic Notation}
	
	Let $I \subseteq \mathbb{R}$ be an interval.  
	A \emph{complex--valued function of a real variable} is a mapping
	$$
	\theta : I \longrightarrow \mathbb{C}.
	$$
	Every such function can be decomposed as
	$$
	\theta(x) = u(x) + i v(x), \qquad u,v : I \to \mathbb{R}.
	$$
	
	\paragraph{Examples.}
	\begin{itemize}
		\item Purely real:  
		$$
		\theta(x) = x^2 = x^2 + i \cdot 0.
		$$
		\item Purely imaginary:  
		$$
		\theta(x) = i e^x = 0 + i e^x.
		$$
		\item General case:  
		$$
		\theta(x) = e^{ix} = \cos x + i \sin x.
		$$
	\end{itemize}
	
	\subsection*{Pedantic Notation for Differential Equations}
	
	A first--order ordinary differential equation is formally written as
	$$
	\theta'(x) = f\bigl(x, \theta(x)\bigr), \qquad x \in I,
	$$
	where
	$$
	f : D \subseteq I \times \mathbb{C} \;\longrightarrow\; \mathbb{C}.
	$$
	
	Here:
	\begin{itemize}
		\item $I \subseteq \mathbb{R}$ is the interval of the independent variable.
		\item $D \subseteq I \times \mathbb{C}$ is the domain of $f$, consisting of all ordered pairs $(x,y)$ with $x \in I$ and $y \in \mathbb{C}$.
		\item $f(x,y) \in \mathbb{C}$ gives the derivative at the point where $\theta(x) = y$.
	\end{itemize}
	
	Thus, $D$ can be visualized as the set of all possible states $(x,y)$ where the equation is defined.
	
	\paragraph{Example.}  
	If $I = \mathbb{R}$ and
	$$
	f(x,y) = xy,
	$$
	then
	$$
	D = \mathbb{R} \times \mathbb{C}, \qquad f : D \to \mathbb{C},
	$$
	and the differential equation reads
	$$
	\theta'(x) = x \, \theta(x).
	$$
	
	\subsection*{Equivalence with Standard Notation}
	
	In practice, the same relation is usually written in the shorthand form
	$$
	y' = f(x,y),
	$$
	with the understanding that
	$$
	y = y(x), \qquad y'(x) = \frac{dy}{dx}, \qquad y(x) \in \mathbb{C}.
	$$
	
	Thus, the pedantic formulation
	$$
	\theta'(x) = f(x,\theta(x))
	$$
	is entirely equivalent to the standard shorthand
	$$
	y' = f(x,y).
	$$
	
	The difference lies only in notation:
	\begin{itemize}
		\item Pedantic form emphasizes $\theta$ as a mapping $I \to \mathbb{C}$.
		\item Standard form suppresses this, using $y$ for the function and its values.
	\end{itemize}
	
	\section*{Pedantic vs. Normal ODE Notation}
	
	\subsection*{Pedantic Formulation}
	
	Let $I \subseteq \mathbb{R}$ be an interval.  
	A first--order ODE is written as
	$$
	\theta'(x) = f\bigl(x,\theta(x)\bigr), \qquad x \in I,
	$$
	where
	$$
	f : D \subseteq I \times \mathbb{C} \;\longrightarrow\; \mathbb{C}.
	$$
	
	Here:
	\begin{itemize}
		\item $I$ is the interval of the independent variable $x$.
		\item $D$ is a subset of $I \times \mathbb{C}$, consisting of pairs $(x,y)$ where the formula for $f(x,y)$ is defined.
		\item $\theta : I \to \mathbb{C}$ is the unknown function to be determined.
	\end{itemize}
	
	\paragraph{Example (Pedantic).}  
	$$
	I = \mathbb{R}, \quad 
	D = I \times \mathbb{C}, \quad
	f(x,y) = x y,
	$$
	so the equation is
	$$
	\theta'(x) = x \, \theta(x).
	$$
	
	\subsection*{Normal Shorthand Formulation}
	
	The same ODE is usually written as
	$$
	y' = f(x,y),
	$$
	with the understanding that
	$$
	y = y(x), \qquad y'(x) = \frac{dy}{dx}.
	$$
	
	In this shorthand,
	\begin{itemize}
		\item $f$ is viewed simply as a function of two variables,
		$$
		f : D \subseteq \mathbb{R}^2 \;\longrightarrow\; \mathbb{R}
		$$
		in the purely real case, or
		$$
		f : D \subseteq \mathbb{R} \times \mathbb{C} \;\longrightarrow\; \mathbb{C}
		$$
		in the complex case.
		\item $D$ is the domain of definition of $f(x,y)$, i.e. the set of all pairs $(x,y)$ where $f$ is defined.
	\end{itemize}
	
	\paragraph{Example (Shorthand).}  
	$$
	y' = xy, \qquad f(x,y) = xy,
	$$
	with
	$$
	D = \mathbb{R}^2.
	$$
	
	\subsection*{Equivalence}
	
	Thus:
	\begin{align*}
		\text{Pedantic: } & \theta'(x) = f\bigl(x,\theta(x)\bigr), 
		\quad f : D \subseteq I \times \mathbb{C} \to \mathbb{C}, \\[6pt]
		\text{Normal: } & y' = f(x,y), 
		\quad f : D \subseteq \mathbb{R}^2 \to \mathbb{R} \; \text{ or } \; f : D \subseteq \mathbb{R} \times \mathbb{C} \to \mathbb{C}.
	\end{align*}
	
	Both notations describe the same situation.  
	The difference is that the pedantic version highlights $\theta$ as a function $I \to \mathbb{C}$ and makes $D \subseteq I \times \mathbb{C}$ explicit, while the shorthand suppresses these details.
	
	\subsection*{Role of the Domain $D$}
	
	In the differential equation
	$$
	\theta'(x) = f\bigl(x,\theta(x)\bigr),
	$$
	the right--hand side $f$ depends on two variables: the independent variable $x \in I$ and the dependent variable $y = \theta(x) \in \mathbb{C}$.  
	Therefore we write
	$$
	f : D \subseteq I \times \mathbb{C} \;\longrightarrow\; \mathbb{C},
	$$
	where $D$ is the set of all pairs $(x,y)$ for which $f(x,y)$ is defined.  
	In this way, $D$ specifies the valid region of the $(x,y)$--plane on which the ODE is meaningful.
	
	\paragraph{Example.}
	If
	$$
	f(x,y) = \frac{y}{x},
	$$
	then $x$ cannot be zero. Hence the domain is
	$$
	D = \{(x,y) \in \mathbb{R} \times \mathbb{C} : x \neq 0 \}.
	$$
	The ODE reads
	$$
	\theta'(x) = \frac{\theta(x)}{x}, \qquad x \neq 0.
	$$
	Thus, $D$ excludes the vertical line $x=0$ in the $(x,y)$--plane, and the differential equation is only defined on regions where $x \neq 0$.
	\subsection*{Examples in Both Notations}
	
	\paragraph{Example 1: Linear Growth.}
	
	\emph{Pedantic notation:}
	$$
	\theta'(x) = f\bigl(x,\theta(x)\bigr), \qquad 
	f(x,y) = y, \quad D = \mathbb{R} \times \mathbb{C}.
	$$
	That is,
	$$
	\theta'(x) = \theta(x).
	$$
	The solution is
	$$
	\theta(x) = C e^x, \qquad C \in \mathbb{C}.
	$$
	
	\emph{Normal shorthand:}
	$$
	y' = f(x,y) = y, \qquad D = \mathbb{R}^2,
	$$
	with the same solution
	$$
	y(x) = C e^x.
	$$
	
	\bigskip
	
	\paragraph{Example 2: Division by $x$.}
	
	\emph{Pedantic notation:}
	$$
	\theta'(x) = f\bigl(x,\theta(x)\bigr), \qquad 
	f(x,y) = \frac{y}{x}, \quad D = \{ (x,y) \in \mathbb{R} \times \mathbb{C} : x \neq 0 \}.
	$$
	That is,
	$$
	\theta'(x) = \frac{\theta(x)}{x}, \qquad x \neq 0.
	$$
	The solution is
	$$
	\theta(x) = Cx, \qquad C \in \mathbb{C}.
	$$
	
	\emph{Normal shorthand:}
	$$
	y' = f(x,y) = \frac{y}{x}, \qquad D = \{(x,y) \in \mathbb{R}^2 : x \neq 0\},
	$$
	with the same solution
	$$
	y(x) = Cx.
	$$
	
\end{document}
